\section{部署思路与治理}
\subsection{模型出口策略}
Weston 强调，self-rewarding 的出口不应直接暴露 raw verified response：需要再经过 alignment 层（meta judge + human checker）确认后才输出。部署策略通常包含：
\begin{itemize}
    \item Deployment filters：只有通过 verified response 与 judge 的一致性检查，才可进入 downstream API。
    \item Conservative fallback：当 judge confidence 低于阈值时，自动降级到 baseline RLHF 模型。
    \item Update schedule：每日/每周生成 self-rewarding 版本，并用 offline dataset 回测 hallucination。
\end{itemize}

\begin{warningbox}{部署风险提示}
如果 self-rewarding 模型在生产中持续 hallucination，则即便 verified response 被 judge 采纳，也可能在对话中产生 misleading explanations。需要 multi-agent cross-check 或者 nightly calibration 来防止 reward model drift。
\end{warningbox}

\subsection{治理与审计}
治理层面要记录：1）每个 verified response 的 judge log；2）reward score 与 compute cost；3）prompt template 与 Self-Instruct seed 的版本。借助这些 log，可以为 future audits 提供透明度，并支撑 `Why did the model take this reasoning path?` 的问询。

\begin{table}[H]
\centering
\begin{tabular}{p{5cm}p{10cm}}
\toprule
治理维度 & 做法 \\
\midrule
审计 trail & 保存 verified response、judge 结果、reward score 以供后续重放。 \\
Prompt provenance & 为每个 Self-Instruct prompt 戳上 seed ID + generation timestamp。 \\
Compute budget governance & 设定 cross-check 最大 token，超出即触发 alert。 \\
\bottomrule
\end{tabular}
\caption{治理 layer 使 self-improvement 可审计。}
\end{table}

\subsection{本章小结}
部署前的过滤策略及治理日志确保 verified response 不会直接造成误导，并让 team 能随时回答 ``How did this reasoning chain earn its reward''。

